\chapter{Chapter05 — DNF Sparse Polynomial Threshold}
%%% generated by the blueprint pipeline from FABL.Chapter05.DNFSparsePolynomialThreshold
\section{Overview}

Declarations of the Lean module \texttt{FABL.Chapter05.DNFSparsePolynomialThreshold}.

\section{Declarations}

\begin{definition}[True indicator of a DNF term]
\lean{FABL.DNFTerm.trueIndicator}
\statementsource{Analysis-of-Boolean-Functions-by-Ryan-ODonnell}{pdf-ffbe9c7b3306-p093-b004}
\label{FABL.DNFTerm.trueIndicator}
\leanok
\uses{FABL.DNFTerm, FABL.DNFTerm.eval}
For a DNF term $T$ and a sign vector $x$, the true indicator of $T$ at $x$ is $1$ if the
term evaluates to $-1$, and $0$ otherwise.
\end{definition}

\begin{theorem}[True indicator characterization]
\lean{FABL.DNFTerm.trueIndicator_eq_one_iff}
\label{FABL.DNFTerm.trueIndicator_eq_one_iff}
\leanok
\uses{FABL.DNFTerm, FABL.DNFTerm.eval, FABL.DNFTerm.trueIndicator}
\difficulty{1}
For a DNF term $T$ on $n$ coordinates and an input $x\in\{\pm1\}^n$, its true indicator
equals $1$ if and only if the term evaluates to $-1$ at $x$.
\end{theorem}

\begin{theorem}[Nonnegativity of the true indicator]
\lean{FABL.DNFTerm.trueIndicator_nonneg}
\label{FABL.DNFTerm.trueIndicator_nonneg}
\leanok
\uses{FABL.DNFTerm, FABL.DNFTerm.eval, FABL.DNFTerm.trueIndicator}
\difficulty{4}
For every DNF term $T$ on $n$ coordinates and every input $x\in\{\pm1\}^n$, the true
indicator of $T$ at $x$ is nonnegative.
\end{theorem}

\begin{definition}[Binary base point of a DNF term]
\lean{FABL.DNFTerm.binaryBasePoint}
\label{FABL.DNFTerm.binaryBasePoint}
\leanok
\uses{FABL.DNFTerm, FABL.DNFTerm.requiredAt, FABL.DNFTerm.support, FABL.binarySignEquiv}
Let $T$ be a DNF term on $n$ coordinates, with distinct literal indices, and let its
support be the set of indices appearing in its literals. Define a point of the binary
cube by assigning to each coordinate in the support the binary value corresponding to
its required sign, and assigning $0$ to every coordinate outside the support.
\end{definition}

\begin{theorem}[Term indicator under binary sign encoding]
\lean{FABL.DNFTerm.trueIndicator_binaryCubeSignEquiv}
\label{FABL.DNFTerm.trueIndicator_binaryCubeSignEquiv}
\leanok
\uses{FABL.DNFTerm, FABL.DNFTerm.binaryBasePoint, FABL.DNFTerm.eval, FABL.DNFTerm.eval_eq_neg_one_iff_requiredAt, FABL.DNFTerm.requiredAt, FABL.DNFTerm.support, FABL.DNFTerm.trueIndicator, FABL.F₂DecisionTree.coordinateZeroSubspace, FABL.F₂DecisionTree.mem_coordinateZeroSubspace_iff, FABL.binaryAffineSubspace, FABL.binaryCubeSignEquiv, FABL.binaryCubeSignEquiv_apply, FABL.binarySignEquiv, FABL.mem_binaryAffineSubspace_iff_add_mem, FABL.setIndicator, FABL.signEncode}
\difficulty{6}
Let $T$ be a DNF term on $n$ coordinates. Under the coordinatewise bijection from binary
vectors to sign vectors, the indicator of the event that $T$ evaluates to $-1$ equals
the indicator of the affine subspace consisting of the satisfying binary base point of
$T$ plus the subspace of vectors that vanish on the support of $T$.
\end{theorem}

\begin{theorem}[Fourier one norm of a term indicator]
\lean{FABL.DNFTerm.fourierOneNorm_trueIndicator}
\label{FABL.DNFTerm.fourierOneNorm_trueIndicator}
\leanok
\uses{FABL.DNFTerm, FABL.DNFTerm.binaryBasePoint, FABL.DNFTerm.support, FABL.DNFTerm.trueIndicator, FABL.DNFTerm.trueIndicator_binaryCubeSignEquiv, FABL.F₂DecisionTree.coordinateZeroSubspace, FABL.fourierOneNorm, FABL.fourierOneNorm_eq_spectralPNorm_one, FABL.spectralPNorm_one_setIndicator_binaryAffineSubspace}
\difficulty{3}
For every conjunction $T$ of literals on $n$ coordinates, the Fourier one norm of its
indicator function—the function equal to $1$ where $T$ evaluates to $-1$ and $0$
elsewhere—is $1$.
\end{theorem}

\begin{theorem}[Fourier coefficient of a finite sum]
\lean{FABL.fourierCoeff_sum}
\label{FABL.fourierCoeff_sum}
\leanok
\uses{FABL.fourierCoeff, FABL.monomial}
\difficulty{3}
Let $n$ be a nonnegative integer, let $\iota$ be a finite index set, let $f_i$ be a
real-valued function on the $n$-dimensional Boolean cube for each $i\in\iota$, and let
$S\subseteq[n]$. The Fourier coefficient at $S$ of $\sum_{i\in\iota} f_i$ is
$\sum_{i\in\iota}\widehat{f_i}(S)$.
\end{theorem}

\begin{theorem}[Fourier norm of a finite sum]
\lean{FABL.fourierOneNorm_sum_le}
\label{FABL.fourierOneNorm_sum_le}
\leanok
\uses{FABL.fourierCoeff, FABL.fourierCoeff_sum, FABL.fourierOneNorm}
\difficulty{3}
Let $n$ be a nonnegative integer, let $I$ be a finite set, and let
$f_i:\\{-1,1\\}^n\\to\\mathbb R$ for each $i\\in I$. The Fourier one-norm of their
pointwise sum is at most the sum of their Fourier one-norms.
\end{theorem}

\begin{definition}[Satisfied term count]
\lean{FABL.DNFFormula.satisfiedTermCount}
\statementsource{Analysis-of-Boolean-Functions-by-Ryan-ODonnell}{pdf-ffbe9c7b3306-p093-b004}
\label{FABL.DNFFormula.satisfiedTermCount}
\leanok
\uses{FABL.DNFFormula, FABL.DNFTerm.trueIndicator}
For a DNF formula $\varphi$ and a sign vector $x$, the satisfied term count is the sum,
over all terms of $\varphi$, of the indicator that the term evaluates to $-1$.
\end{definition}

\begin{theorem}[Zero count when no term is satisfied]
\lean{FABL.DNFFormula.satisfiedTermCount_eq_zero_of_no_satisfied_term}
\label{FABL.DNFFormula.satisfiedTermCount_eq_zero_of_no_satisfied_term}
\leanok
\uses{FABL.DNFFormula, FABL.DNFFormula.satisfiedTermCount, FABL.DNFTerm.eval, FABL.DNFTerm.trueIndicator}
\difficulty{4}
Let $n\in\mathbb N$, let $\varphi$ be a disjunctive normal form on $n$ coordinates, and
let $x\in\{-1,1\}^n$. If no term of $\varphi$ evaluates to $-1$ at $x$, then the number
of terms evaluating to $-1$ at $x$ is zero.
\end{theorem}

\begin{theorem}[A satisfied term gives a positive term count]
\lean{FABL.DNFFormula.one_le_satisfiedTermCount_of_satisfied_term}
\label{FABL.DNFFormula.one_le_satisfiedTermCount_of_satisfied_term}
\leanok
\uses{FABL.DNFFormula, FABL.DNFFormula.satisfiedTermCount, FABL.DNFTerm, FABL.DNFTerm.eval, FABL.DNFTerm.trueIndicator, FABL.DNFTerm.trueIndicator_eq_one_iff, FABL.DNFTerm.trueIndicator_nonneg}
\difficulty{4}
Let $n$ be a natural number, let $\varphi$ be a DNF formula on the sign cube
$\{-1,1\}^n$, and let $x$ be an input in that cube. If a term $T$ of $\varphi$ is
satisfied by $x$, then the number of terms of $\varphi$ satisfied by $x$ is at least
$1$.
\end{theorem}

\begin{theorem}[Fourier norm of the satisfied-term count]
\lean{FABL.DNFFormula.fourierOneNorm_satisfiedTermCount_le}
\label{FABL.DNFFormula.fourierOneNorm_satisfiedTermCount_le}
\leanok
\uses{FABL.DNFFormula, FABL.DNFFormula.satisfiedTermCount, FABL.DNFFormula.size, FABL.DNFTerm.fourierOneNorm_trueIndicator, FABL.DNFTerm.trueIndicator, FABL.fourierOneNorm, FABL.fourierOneNorm_sum_le}
\difficulty{4}
Let $n$ be a natural number and let $\varphi$ be a DNF formula on the sign cube
$\{-1,1\}^n$. The Fourier one-norm of the function counting the terms of $\varphi$
satisfied by an input is at most the number of terms in $\varphi$.
\end{theorem}

\begin{theorem}[Threshold representation of DNF evaluation]
\lean{FABL.DNFFormula.toBooleanFunction_eq_thresholdSign_satisfiedTermCount}
\label{FABL.DNFFormula.toBooleanFunction_eq_thresholdSign_satisfiedTermCount}
\leanok
\uses{FABL.DNFFormula, FABL.DNFFormula.eval_eq_neg_one_iff, FABL.DNFFormula.one_le_satisfiedTermCount_of_satisfied_term, FABL.DNFFormula.satisfiedTermCount, FABL.DNFFormula.satisfiedTermCount_eq_zero_of_no_satisfied_term, FABL.DNFFormula.toBooleanFunction, FABL.DNFTerm.eval, FABL.thresholdSign, FABL.thresholdSign_of_neg, FABL.thresholdSign_of_nonneg}
\difficulty{5}
Let $n\in\mathbb{N}$, let $\varphi$ be a DNF formula on $n$ sign coordinates, and let
$x$ be an input. The value of the formula at $x$ equals the threshold sign of one half
minus the number of terms satisfied by $x$.
\end{theorem}

\begin{theorem}[Sparsity after a half-shift]
\lean{FABL.polynomialSparsity_half_sub_le}
\label{FABL.polynomialSparsity_half_sub_le}
\leanok
\uses{FABL.expect_monomial, FABL.fourierCoeff, FABL.fourierCoeff_sub, FABL.fourierSupport, FABL.mem_fourierSupport, FABL.polynomialSparsity}
\difficulty{5}
For every nonnegative integer $n$ and every real-valued function $q$ on the
$n$-dimensional sign cube, the number of nonzero Fourier coefficients of the function
$x\mapsto \tfrac12-q(x)$ is at most one more than the number of nonzero Fourier
coefficients of $q$.
\end{theorem}

\begin{theorem}[Quadratic sparsity threshold representation]
\lean{FABL.DNFFormula.exists_polynomialThresholdRepresentation_sparsity_le_quadratic}
\proofsource{Analysis-of-Boolean-Functions-by-Ryan-ODonnell}{
  pdf-ffbe9c7b3306-p115-b005,
  pdf-ffbe9c7b3306-p117-b004,
  pdf-ffbe9c7b3306-p117-b005,
  pdf-ffbe9c7b3306-p135-b007
}
\label{FABL.DNFFormula.exists_polynomialThresholdRepresentation_sparsity_le_quadratic}
\leanok
\uses{FABL.DNFFormula, FABL.DNFFormula.eval, FABL.DNFFormula.fourierOneNorm_satisfiedTermCount_le, FABL.DNFFormula.one_le_satisfiedTermCount_of_satisfied_term, FABL.DNFFormula.satisfiedTermCount, FABL.DNFFormula.satisfiedTermCount_eq_zero_of_no_satisfied_term, FABL.DNFFormula.size, FABL.DNFFormula.toBooleanFunction, FABL.DNFFormula.toBooleanFunction_eq_thresholdSign_satisfiedTermCount, FABL.DNFTerm.eval, FABL.IsPolynomialThresholdRepresentation, FABL.exists_sparsePolynomial_uniformApproximation, FABL.fourierCoeff, FABL.fourierOneNorm, FABL.fourierOneNorm_nonneg, FABL.fourierSupport, FABL.polynomialSparsity, FABL.polynomialSparsity_half_sub_le, FABL.thresholdSign, FABL.thresholdSign_of_neg, FABL.thresholdSign_of_nonneg}
\difficulty{6}
Let $n$ be a positive integer and let $\varphi$ be a DNF formula on $n$ sign-valued
Boolean variables. There is a real-valued polynomial on the Boolean cube whose threshold
sign agrees with the function computed by $\varphi$ at every input, and whose Fourier
sparsity is at most $17n$ times the square of the number of terms of $\varphi$.
\end{theorem}

\begin{theorem}[Exists polynomial threshold representation of has dnf size le sparsity le quadratic]
\lean{FABL.exists_polynomialThresholdRepresentation_of_hasDNFSizeLE_sparsity_le_quadratic}
\statementsource{Analysis-of-Boolean-Functions-by-Ryan-ODonnell}{
  pdf-ffbe9c7b3306-p135-b007,
  pdf-ffbe9c7b3306-p093-b004
}
\label{FABL.exists_polynomialThresholdRepresentation_of_hasDNFSizeLE_sparsity_le_quadratic}
\leanok
\uses{FABL.BooleanFunction, FABL.DNFFormula.exists_polynomialThresholdRepresentation_sparsity_le_quadratic, FABL.HasDNFSizeLE, FABL.IsPolynomialThresholdRepresentation, FABL.polynomialSparsity}
\difficulty{3}
% dataset-pipeline: stub -- prose pending build_informal_statements.py
Statement pending regeneration.
\end{theorem}

\begin{theorem}[Exists polynomial threshold representation of has dnf size le sparsity le cubic]
\lean{FABL.exists_polynomialThresholdRepresentation_of_hasDNFSizeLE_sparsity_le_cubic}
\statementsource{Analysis-of-Boolean-Functions-by-Ryan-ODonnell}{
  pdf-ffbe9c7b3306-p093-b004,
  pdf-ffbe9c7b3306-p135-b007
}
\label{FABL.exists_polynomialThresholdRepresentation_of_hasDNFSizeLE_sparsity_le_cubic}
\leanok
\uses{FABL.BooleanFunction, FABL.HasDNFSizeLE, FABL.IsPolynomialThresholdRepresentation, FABL.exists_polynomialThresholdRepresentation_of_hasDNFSizeLE_sparsity_le_quadratic, FABL.polynomialSparsity}
\difficulty{4}
% dataset-pipeline: stub -- prose pending build_informal_statements.py
Statement pending regeneration.
\end{theorem}
